% =============================================================================
% Descrição do projeto. Itens 2 a 9 do Anexo I, na ordem do Anexo.
%   1.1 Introdução                     (item 2)
%   1.2 Objetivos                      (item 3)
%   1.3 Público alvo                   (item 4)
%   1.4 Metodologia de desenvolvimento (item 5)
%   1.5 Recursos necessários           (item 6)
%   1.6 Resultados fundamentais        (item 7)
%   1.7 Contribuição / formação        (item 8)
%   1.8 Cronograma                     (item 9)
% O item 1 (título) está na folha de rosto; o item 10, em REFERÊNCIAS.
%
% Convenções de escrita adotadas neste arquivo: sem travessões, sem negrito e sem
% realce no corpo do texto. Itálico apenas em palavras estrangeiras, como pede a
% norma. Vocabulário simples, formal e objetivo.
% =============================================================================

\chapter{Descrição do projeto}
\label{cap:descricao}

\section{Introdução}
\label{sec:introducao}

Interfaces digitais atuais, como PCI Express Gen 6, DDR5 e Ethernet de 112\,Gb/s,
impõem tempos de transição de dezenas de picossegundos, o que leva conteúdo
espectral relevante para dezenas de gigahertz. Nesse regime, trilhas comportam-se
como linhas de transmissão, pares de planos como cavidades ressonantes e
descontinuidades geométricas como antenas não intencionais
\citep{johnson1993,ott2009}. A rede de distribuição de energia (\textit{power
delivery network}, PDN) ocupa posição central: os modos próprios da cavidade
formada pelos planos produzem picos na curva de impedância $Z(f)$, que amplificam o
ruído de comutação e o convertem em emissão irradiada pelas bordas da placa e pelos
cabos a ela conectados \citep{bogatin2010}.

O projetista, porém, não formula sua pergunta em termos de impedância, e sim de
forma normativa: este empilhamento tem margem contra a curva-limite de emissão da
norma que o produto precisa atender, no ambiente em que será usado? Hoje a resposta
chega apenas ao fim do ciclo de desenvolvimento, por ensaio em laboratório
acreditado sobre protótipo. A reprovação nesse ponto implica reprojeto, nova
fabricação e novo agendamento de ensaio, com impacto medido em meses. Antecipar
essa resposta exige avaliar milhares de configurações, e é aí que o custo
computacional se torna alto demais: solvers de onda completa consomem de dezenas de
minutos a horas por configuração, e mesmo métodos baseados em física, uma a duas
ordens de grandeza mais rápidos \citep{schierholz2021,hillebrecht2024}, seguem
incompatíveis com otimização automatizada.

Modelos substitutos deslocam esse custo para uma fase de treinamento única e
entregam predições em milissegundos \citep{swaminathan2020,jiang2020,zhang2021}.
Persiste, contudo, uma limitação estrutural: modelos orientados apenas a dados
reduzem o erro numérico sem garantir consistência física, e redes treinadas sobre
parâmetros de espalhamento produzem com frequência respostas não passivas ou não
causais, inúteis para simulação no domínio do tempo \citep{torun2020}. O problema se
agrava quanto menor a quantidade de dados disponível, pois a rede não tem evidência
suficiente para inferir, apenas deles, comportamentos que a física já descreve em
forma fechada \citep{raissi2019}.

Propõe-se, por isso, uma arquitetura de duas camadas. A primeira é um modelo
substituto informado por física, que prediz $Z(f)$ a partir dos parâmetros do
empilhamento, com restrições eletromagnéticas incorporadas à função de perda. A
segunda é uma camada normativa determinística, que converte $Z(f)$ em uma
estimativa de emissão e a compara com a curva-limite da norma e do ambiente
selecionados, produzindo uma margem em decibéis. É essa margem, e não a impedância,
que serve de função objetivo à otimização. Uma investigação preliminar já realizada
sobre a base mostrou quais restrições podem ser usadas: o regime quase-estático
segue com precisão a lei capacitiva prevista, com inclinação logarítmica medida de
$-1{,}018 \pm 0{,}008$ contra $-1$ teórico, enquanto o uso das frequências de
ressonância como referência, inicialmente previsto como principal, foi refutado.
Formula-se o problema de pesquisa:

\begin{quote}
Em que medida a incorporação de restrições eletromagnéticas conhecidas em forma
fechada à função de perda de um modelo substituto neural melhora a acurácia, a
consistência física e a generalização entre topologias na predição da impedância de
PDNs, em função do volume de dados de treinamento, e em que medida a margem
normativa dela derivada ordena empilhamentos de forma estável o suficiente para
orientar decisões de projeto?
\end{quote}

\section{Objetivos}
\label{sec:objetivos}

\subsection{Objetivo geral}

Desenvolver e validar um modelo substituto informado por física que, acoplado a uma
camada normativa determinística, estime a margem de conformidade eletromagnética da
rede de distribuição de energia de uma placa de circuito impresso multicamadas
contra as curvas-limite de emissão da norma e do ambiente de uso selecionados pelo
projetista, com acurácia e latência compatíveis com a otimização do empilhamento na
fase inicial de projeto.

\subsection{Objetivos específicos}

\begin{alineas}
\item implementar um pipeline incremental de extração da curva de impedância a
partir dos arquivos Touchstone, validado contra propriedades analíticas conhecidas,
e reunir os subconjuntos de PDN da SI/PI-Database, de 1 a 13 cavidades, em uma
representação de atributos independente da topologia;

\item verificar sobre os próprios dados quais relações analíticas da cavidade são
observáveis na grandeza predita, formulando os termos de perda apenas a partir das
que se confirmarem;

\item treinar e avaliar o modelo informado por física contra um conjunto de modelos
de referência, sob protocolo único, medindo a consistência física, o ganho em função
do volume de treinamento e a generalização para topologias não vistas;

\item implementar a camada normativa de estimativa de margem, medindo a estabilidade
do ordenamento de empilhamentos sob variação de suas premissas, e acoplá-la a um
otimizador evolutivo multiobjetivo que maximize a margem e minimize o custo de
fabricação estimado;

\item obter diretrizes de projeto interpretáveis, expressas em decibéis de margem, e
publicar a implementação como software livre reprodutível.
\end{alineas}

\section{Público alvo}
\label{sec:publico}

O público primário é o engenheiro de projeto de \textit{hardware} digital e a equipe
de compatibilidade eletromagnética que o apoia, em organizações cujos produtos
passam por ensaios de certificação. O momento de uso é a fase inicial de definição
do empilhamento, quando número de camadas, espessura de cavidade, dielétrico e
densidade de vias ainda podem ser alterados. O público secundário é a comunidade
acadêmica de integridade de sinal e potência e de aprendizado informado por física,
para a qual se destinam a implementação aberta, o protocolo experimental e as
hipóteses formuladas de modo que possam ser refutadas. Há ainda o público
institucional: o pipeline, a base reunida e a camada normativa ficam disponíveis
para trabalhos futuros no Departamento de Engenharia Elétrica da UFPR.

\section{Metodologia de desenvolvimento}
\label{sec:metodologia}

\subsection{Base de dados}
\label{subsec:harmonizacao}

Utiliza-se a SI/PI-Database, do \textit{Institut für Theoretische Elektrotechnik} da
Universidade Técnica de Hamburgo \citep{schierholz2021,hillebrecht2024}, de acesso
público mediante aceite dos termos de uso. A referência controlada é o subconjunto
\textit{6-Layer PCB based PDN with Two Via Arrays}, amostrado por hipercubo latino e
já caracterizado por inspeção direta dos arquivos: 985 configurações, 36 portas, 334
pontos entre 1\,MHz e 1\,GHz e oito graus de liberdade efetivos entre os dezessete
parâmetros tabulados.

O trabalho não se restringe a esse subconjunto. A base oferece catorze subconjuntos
de PDN, com 1 a 13 cavidades e cerca de 135 mil configurações, todos cobrindo 1\,MHz
a 1\,GHz e simulados pela mesma ferramenta baseada em física. Eles são combináveis
porque o alvo é a autoimpedância em uma única porta, de modo que a variação no
número total de portas, de 2 a 68, não altera a dimensão da saída. Exclui-se apenas
o subconjunto de onda completa de 14 camadas, fora da banda e de fidelidade
distinta. A união exige quatro tratamentos. O espaço de parâmetros é diferente em
cada subconjunto, o que leva a adotar atributos derivados independentes da
topologia, como número de cavidades, espessura total, capacitância acumulada,
densidade de vias, passo normalizado e razões geométricas. A porta de interesse
recebe numeração distinta em cada subconjunto, o que exige um adaptador conferido
contra a folha de dados. As variantes Eurocard usam 121 pontos de frequência na
mesma faixa, e são interpoladas para a grade comum. E os 119\,GB comprimidos do
conjunto completo excederiam o armazenamento se descompactados de uma vez, de modo
que cada subconjunto é baixado, reduzido à curva de impedância e tem os arquivos
brutos descartados antes do próximo. Como os desenhos de amostragem diferem, o
experimento controlado de curva de aprendizado fica restrito aos subconjuntos por
hipercubo latino.

O ganho não é só de volume. Com o número de cavidades variando de 1 a 13, essa
grandeza passa a ser um atributo observável, o que permite testar se a capacitância
efetiva cresce com o número de cavidades em paralelo, desvio que a investigação
preliminar detectou e não pôde explicar com um único subconjunto. Confirmada a
relação, o termo quase-estático passa a restringir também a magnitude da resposta. A
variedade de topologias permite ainda validação com topologia retida, o que torna
mensurável a principal limitação de validade externa do trabalho.

\subsection{Modelo informado por física}

O alvo é a curva $y(f) = \log_{10}|Z_{11}(f)|$, obtida da matriz de espalhamento.
Prediz-se a curva completa, e não um valor escalar, porque é ela que alimenta a
camada normativa e porque as restrições físicas atuam sobre sua forma. A função de
perda soma o erro de reconstrução, ponderado por década para compensar a
subamostragem da primeira década, a três termos derivados da física: o desvio da
derivada logarítmica em relação a $-1$ na região quase-estática, formulação que é
invariante a fator multiplicativo e por isso não é afetada pelo desvio de escala
observado; a monotonicidade até o nulo de série, que penaliza oscilações falsas,
comuns em redes treinadas com poucos dados; e a passividade. Os pesos são
hiperparâmetros críticos e serão determinados por otimização bayesiana sobre o
conjunto de validação, com a sensibilidade do resultado reportada.

Os modelos de referência, avaliados sob protocolo idêntico, são regressão
regularizada, floresta aleatória, \textit{gradient boosting}, processos gaussianos e
um perceptron multicamadas \citep{geron2022}. Este último é o mais importante, por
diferir do modelo proposto apenas pela ausência dos termos físicos, o que isola o
efeito da informação física. O conjunto é particionado com semente registrada, o
teste é acessado uma única vez e o ganho é medido em frações crescentes do
treinamento, com significância avaliada por teste não paramétrico pareado. O
critério de aceitação é $R^2 > 0{,}90$ sobre $|Z|_{\max}$ em teste isolado, com
latência de inferência inferior a 1\,ms.

\subsection{Camada normativa de estimativa de margem}
\label{subsec:normativa}

A conversão de $Z(f)$ em margem é determinística, sem aprendizado, e ocorre em três
etapas com premissas declaradas. A excitação não integra a base de dados e é entrada
do projetista: adota-se um modelo de corrente trapezoidal, definido por frequência
de relógio, tempo de transição e amplitude, cujo envelope espectral decai a
$-20$\,dB por década até o joelho em $1/(\pi t_r)$ e a $-40$\,dB por década acima
dele \citep{johnson1993}, do que resulta $V_n(f) = |Z_{11}(f)|\,I(f)$. A radiação é
modelada por expressões fechadas para os dois mecanismos dominantes abaixo de
1\,GHz \citep{ott2009}: a emissão pelas bordas do par de planos e a emissão de modo
comum excitada por cabos. A comparação normativa usa as curvas-limite da CISPR 32,
classes A e B \citep{cispr32}, e das normas genéricas IEC 61000-6-3
\citep{iec610006m3} e IEC 61000-6-4 \citep{iec610006m4}, selecionáveis pelo ambiente
de uso. A margem é $m(f) = L(f) - \hat{E}(f)$ e a medida de desempenho é a menor
margem na faixa.

Duas delimitações são essenciais. Os limites de emissão irradiada começam em
30\,MHz, de modo que a margem é calculada sobre 30\,MHz a 1\,GHz. Além disso, a
camada é analítica e não é calibrada contra medição, pois a base não contém dados de
campo. Sua saída não é uma decisão de certificação, que depende também de gabinete,
cabos, aterramento e \textit{firmware}, e sim um indicador de triagem, válido para
comparar projetos entre si. Essa validade será medida, e não suposta: varre-se um
conjunto de premissas plausíveis e avalia-se a estabilidade do ordenamento pelo
coeficiente de Kendall.

\subsection{Otimização, interpretabilidade e validação}
\label{subsec:validacao}

Validado o conjunto, ele é usado como função de avaliação em um problema
multiobjetivo que maximiza a menor margem e minimiza um custo de fabricação
estimado, por NSGA-II \citep{deb2002} comparado à evolução diferencial multiobjetivo
\citep{coelho2010}, com as soluções extremas da fronteira de Pareto conferidas por
simulação de referência. A interpretabilidade emprega valores de Shapley
\citep{lundberg2017} e regressão simbólica \citep{tessaro2026} para obter expressões
algébricas que relacionem parâmetros de empilhamento à margem, com faixas de
validade declaradas.

Sob coorientação do Prof.\ Dr.\ Bruno Pohlot Ricobom, prevê-se ainda um estudo de
caso com o escâner de campo próximo desenvolvido pelo autor em iniciação científica:
fabricar duas placas com empilhamentos contrastantes e verificar se o ordenamento de
emissão medido corresponde ao predito. É o único ponto em que a camada normativa é
comparada com medição. Trata-se de verificação de consistência entre ordenamentos, e
não de validação metrológica. Dada a incerteza ainda não caracterizada do
instrumento, nenhuma conclusão quantitativa será extraída desta etapa, que permanece
fora do caminho crítico.

\section{Recursos necessários}
\label{sec:recursos}

A execução é conduzida pelo autor sob orientação do Prof.\ Dr.\ Leandro dos Santos
Coelho, responsável pelas frentes de modelagem substituta, otimização e regressão
simbólica, e sob coorientação do Prof.\ Dr.\ Bruno Pohlot Ricobom, na frente de
instrumentação e medição. Estima-se dedicação de dez horas semanais ao longo de dois
semestres letivos, com reuniões quinzenais de acompanhamento.

Não há necessidade de aquisição de equipamentos ou licenças. Basta um computador
pessoal com 8\,GB de memória e cerca de 60\,GB livres em disco, o que é suficiente
mesmo com todos os subconjuntos, dado o pré-processamento incremental. A cadeia de
software é integralmente livre: numpy, pandas e scipy para manipulação numérica;
scikit-rf para os arquivos Touchstone; scikit-learn e LightGBM para os modelos de
referência; PyTorch para o modelo informado por física; Optuna e pymoo para
otimização; SHAP e PySR para interpretabilidade; e MLflow, pytest, Hydra, Git e
\LaTeX{} para registro, testes, configuração, versionamento e redação. O treinamento
roda em CPU, com o Google Colab gratuito como alternativa. Não há custos previstos no
escopo principal. Um custo de R\$\,300 a R\$\,600 incorreria apenas na fabricação das
placas de teste da validação complementar.

\section{Resultados fundamentais a serem atingidos}
\label{sec:resultados}

\begin{alineas}
\item um modelo substituto informado por física com $R^2 > 0{,}90$ sobre
$|Z|_{\max}$ em teste isolado e latência inferior a 1\,ms, e um conjunto de dados
reunido a partir da base original, reprodutível, cobrindo os subconjuntos de PDN de
1 a 13 cavidades sob representação independente da topologia;

\item a medida do efeito da informação física sobre o erro de generalização, em
função do volume de treinamento e para topologias não vistas, com significância
estatística avaliada;

\item uma camada normativa que traduz a impedância predita em margem em decibéis
contra a CISPR 32 e as normas IEC 61000-6-3 e 6-4, com a estabilidade do ordenamento
medida, e uma fronteira de Pareto de empilhamentos que equilibram margem e custo,
com expressões analíticas interpretáveis e faixas de validade declaradas;

\item uma implementação de código aberto, reprodutível e documentada, e, como
desdobramento pretendido, um manuscrito submetido a evento da área.
\end{alineas}

\section{Contribuição esperada para a ênfase e importância para a formação do
autor}
\label{sec:contribuicao}

A ênfase em sistemas eletrônicos embarcados forma o profissional que projeta o
\textit{hardware} de um produto e responde por seu funcionamento em campo, e é nesse
ponto que o trabalho atua. O sistema embarcado é onde a decisão de empilhamento, o
ruído de comutação e a aprovação em ensaio deixam de ser assuntos separados: a mesma
escolha de espessura de cavidade que reduz o custo da placa determina a impedância
vista pelos componentes e, por consequência, a emissão medida no laboratório. A
contribuição está em tratar essa camada física, em geral repassada a especialistas
ou descoberta tarde por reprovação em ensaio, como objeto explícito de projeto, e em
expressar seu resultado na linguagem que o projetista usa. Há também integração
curricular, pois o problema exige ao mesmo tempo eletromagnetismo aplicado, teoria de
circuitos e linhas de transmissão, análise de sinais no domínio da frequência, normas
de compatibilidade eletromagnética e métodos numéricos.

Para a formação do autor, o trabalho desenvolve três competências. A primeira é
metodológica: conduzir investigação empírica com hipóteses que podem ser refutadas e
com protocolo fixado antes da observação dos resultados. A segunda é técnica: aplicar
aprendizado de máquina a um problema cuja física é parcialmente conhecida em forma
fechada, situação em que cabe ao engenheiro decidir o que o modelo deve aprender e o
que já se sabe. A terceira é de prática profissional: distinguir o que o modelo
sustenta do que não sustenta, distinção que, na camada normativa, separa um
instrumento de triagem de uma afirmação de conformidade indevida.

\section{Cronograma}
\label{sec:cronograma}

O trabalho distribui-se por dois semestres letivos, conforme a
Tabela~\ref{tab:crono}. O período 2026/2 é dedicado à fundamentação e à
infraestrutura de dados, e o período 2027/1 ao modelo informado por física, à camada
normativa e à otimização.

\begin{table}[htb]
\centering
\caption{Cronograma de execução}
\label{tab:crono}
\footnotesize
\begin{tabular}{@{}p{0.38\linewidth}ccccc@{\hspace{1.1em}}ccccc@{}}
\toprule
& \multicolumn{5}{c}{TCC I (2026/2)} &
  \multicolumn{5}{c}{TCC II (2027/1)} \\
\cmidrule(r){2-6}\cmidrule(l){7-11}
Atividade & ago & set & out & nov & dez & mar & abr & mai & jun & jul \\
\midrule
Revisão bibliográfica                        & \mes & \mes & \mes &      &      &      &      &      &      &      \\
Pipeline de extração e validação da base     & \mes & \mes & \mes &      &      &      &      &      &      &      \\
Reunião dos subconjuntos e das normas        &      & \mes & \mes & \mes &      &      &      &      &      &      \\
Modelos de referência e protocolo            &      &      & \mes & \mes &      &      &      &      &      &      \\
Entrega e apresentação de TCC I              &      &      &      &      & \mes &      &      &      &      &      \\
Termos físicos e calibração dos pesos        &      &      &      &      &      & \mes & \mes & \mes &      &      \\
Camada normativa e estabilidade da ordenação &      &      &      &      &      &      & \mes & \mes &      &      \\
Otimização, interpretabilidade e validação   &      &      &      &      &      &      &      & \mes & \mes &      \\
Redação                                      &      &      & \mes & \mes & \mes &      &      & \mes & \mes & \mes \\
Revisão final e defesa                       &      &      &      &      &      &      &      &      &      & \mes \\
\bottomrule
\end{tabular}
\vspace{0.3em}
\\{\footnotesize FONTE: O autor (2026).}
\end{table}

Quatro marcos orientam a continuidade: o pipeline validado contra propriedades
analíticas conhecidas, em setembro de 2026; quatro subconjuntos reunidos e o melhor
modelo de referência em $R^2 > 0{,}80$, em novembro de 2026; o modelo informado por
física superando a ablação com significância estatística, em abril de 2027; e a
fronteira de Pareto conferida por simulação, em junho de 2027. O terceiro marco não
condiciona a validade do trabalho, pois a hipótese pode ser refutada, e sua refutação
é um resultado que também merece publicação.
